\PassOptionsToPackage{hyphens}{url}
\documentclass[preprint,12pt]{elsarticle}
\journal{Computer Methods and Programs in Biomedicine}
\usepackage[T1]{fontenc}
\usepackage{graphicx}
\usepackage{amsmath,amssymb}
\usepackage{booktabs}
\usepackage{array}
\usepackage{algorithm}
\usepackage{algpseudocode}
\usepackage{tikz}
\usepackage{longtable}
\usepackage{lineno}
\usepackage{tabularx}
\usepackage{ragged2e}
\usepackage{adjustbox}
\usepackage{makecell}
\usepackage{pdflscape}
\usepackage[table]{xcolor}
\usepackage{threeparttable}
\usepackage{microtype}
\usepackage{caption}
\captionsetup{font=small,labelfont=bf}
\emergencystretch=2em
\setlength{\LTcapwidth}{\textwidth}
\providecommand{\botrule}{\bottomrule}
\renewcommand{\arraystretch}{1.15}
\usetikzlibrary{arrows.meta,calc,positioning,shapes.geometric,shapes.misc,decorations.pathmorphing}
\newcolumntype{Y}{>{\RaggedRight\arraybackslash}X}
\newcolumntype{C}[1]{>{\Centering\arraybackslash}p{#1}}

% Colors
% IBM Carbon-inspired pastel palette for TikZ Figure 1
\definecolor{ibmBlueLight}{HTML}{D0E2FF}
\definecolor{ibmBlue}{HTML}{78A9FF}
\definecolor{ibmPurpleLight}{HTML}{E8DAFF}
\definecolor{ibmPurple}{HTML}{BE95FF}
\definecolor{ibmCyanLight}{HTML}{BAE6FF}
\definecolor{ibmCyan}{HTML}{33B1FF}
\definecolor{ibmOrangeLight}{HTML}{FFD6A8}
\definecolor{ibmOrange}{HTML}{FF832B}
\definecolor{ibmRedLight}{HTML}{FFD7D9}
\definecolor{ibmRed}{HTML}{FF8389}
\definecolor{ibmGrayLight}{HTML}{F4F4F4}
\definecolor{ibmGray}{HTML}{8D8D8D}
\definecolor{ibmText}{HTML}{262626}
\definecolor{ibmLine}{HTML}{525252}


\usepackage{url}
\usepackage[hidelinks]{hyperref}

\begin{document}

\begin{frontmatter}

\title{A Governance-Aware Rule-Based Computational Framework for Biospecimen Qualification in Biobank Information Systems}

\author[inst1]{Denise Ribeiro\corref{cor1}}
\ead{denise.ribeiro@mail.usf.edu.br}
\author[inst1,inst2]{Fabio Andrijauskas}
\ead{fabio.andrijauskas@usf.edu.br}
\author[inst1]{Lucas Miguel de Carvalho}
\ead{lucas.miguel@usf.edu.br}
\author[inst1]{Vicente Idalberto Becerra Sablon}
\ead{vicente.sablon@usf.edu.br}
\cortext[cor1]{Corresponding author.}

\affiliation[inst1]{
    organization={São Francisco University},
    city={Bragança Paulista},
    postcode={12916-900},
    state={SP},
    country={Brazil}
}

\affiliation[inst2]{
    organization={San Diego Supercomputer Center},
    city={San Diego},
    postcode={10587},
    state={CA},
    country={USA}
}

%% Research highlights
%\begin{highlights}
%\item We propose a governance-aware computational framework for biospecimen qualification.
%\item Derived penalties are separated from non-compensable governance failure.
%\item Synthetic validation achieved high deterministic calibration across 6000 profiles.
%\item Threshold and perturbation analyses demonstrated robust grade stability.
%\item Runtime benchmarking showed scalable execution up to 5 million biospecimens.
%\end{highlights}

\begin{abstract}
\textbf{Background and Objective:} The clinical and molecular integrity of biospecimens stored in biobanks is strongly influenced by 
pre-analytical conditions. Although international standards and best-practice frameworks support biospecimen annotation, metadata 
interoperability, traceability, and governance, they do not directly provide a transparent specimen-level operational quality score. This study 
presents the Ribeiro Classification Score, a governance-aware heuristic framework that transforms Standard PREanalytical Code-derived 
pre-analytical evidence into a standardized operational quality score and maps governance-admissible biospecimens to a five-level A--E 
certification scale.

\textbf{Methods:} The framework was implemented as a deterministic rule-based computational pipeline for biospecimen qualification. Standard 
PREanalytical Code-derived descriptors defined matrix-specific pre-analytical axes, whereas international standards and guidelines informed 
metadata representation, governance admissibility, monitoring, traceability, and documentation requirements. The Ribeiro Classification Score 
starts from a reference value of 100 and subtracts weighted operational penalties associated with biological degradation conditions. A 
non-compensable governance gate independently evaluates admissibility criteria and may directly route biospecimens to Grade E regardless of the 
numerical score. Validation included synthetic-profile analysis, logical consistency checks, sensitivity and threshold analyses, weight-
perturbation experiments, and runtime benchmarking.

\textbf{Results:} All seven property-based logical checks passed successfully. Synthetic validation included 6000 fluid and solid biospecimen 
profiles and yielded an internal calibration index of 0.988833, with 5933 matches and 67 threshold-boundary mismatches between expected and 
assigned grades. Governance-failure profiles were consistently routed to Grade E. Weight-perturbation experiments demonstrated minimal grade 
instability, and runtime benchmarking across 10,000 to 5 million biospecimens showed near-linear scalability (coefficient of determination = 
0.998758; adjusted coefficient of determination = 0.998721).

\textbf{Conclusions:} The Ribeiro Classification Score provides an interpretable and auditable certification framework that links standards-based 
biospecimen evidence to operational quality governance. By explicitly separating governance admissibility failure from pre-analytical penalty 
accumulation, the framework supports reproducible, transparent, and computationally scalable biospecimen qualification for large-scale biobank 
environments.
\end{abstract}

\begin{keyword}
Biobanking; Biospecimen quality; Pre-analytical variability; Metadata interoperability; Biomedical decision support
\end{keyword}

\end{frontmatter}
\linenumbers

\section{Introduction}\label{sec:introduction}

Biobanks are strategic infrastructures for translational research, precision medicine, and long-term biospecimen 
governance \cite{malsagova2020,muller2020}. In this context, the scientific value of a specimen depends not only on 
its biological origin, but also on the quality of its pre-analytical history. This phase is widely recognized as one 
of the most vulnerable components of laboratory practice and can substantially influence molecular integrity, 
analytical reproducibility, and downstream interpretability \cite{plebani2006,campbell2018}.

To reduce pre-analytical ambiguity and strengthen biospecimen governance, international standards and best-practice 
frameworks recommend the use of controlled vocabularies, structured metadata, traceability mechanisms, and quality 
management requirements. SPREC (Standard PREanalytical Code) provides a formal codification of the pre-analytical 
trajectory of fluid and solid biospecimens \cite{betsou2010,lehmann2012}, whereas MIABIS (Minimum Information About 
Biobank data Sharing) supports interoperable metadata representation and exchange across biobank environments 
\cite{norlin2012,merinomartinez2016}. 

In parallel, ISBER Best Practices and ISO 20387 provide guidance and formal requirements for quality management, traceability, biospecimen 
handling, and operational governance in biobanking \cite{campbell2018,iso20387}. However, they do not, by themselves, derive a transparent 
specimen-level operational certification score. This creates a methodological gap between standards-based biospecimen documentation and the 
operational need to convert pre-analytical and governance evidence into reproducible local decision support.

This gap becomes especially relevant in biobank software environments that need to reconcile traceability, semantic interoperability, technical 
governance, and local decision support at the same time \cite{norlin2012,merinomartinez2016,campbell2018,iso20387}. A system may preserve metadata 
correctly and still lack a transparent mechanism for deriving how a specimen should be certified according to its accumulated pre-analytical 
burden and governance admissibility status. In this sense, the problem is not the absence of descriptive standards, but the absence of an explicit 
decision layer capable of converting standards-based evidence into a reproducible local quality status. Current biobank infrastructures frequently 
preserve metadata while lacking reproducible computational decision mechanisms for automated specimen-level qualification.

To address this problem, the proposed framework adopts a deterministic Rule-Based Expert System, in which domain knowledge is encoded as explicit 
logical rules rather than inferred through opaque probabilistic training \cite{shortliffe2014,russell2021}. In healthcare and laboratory 
governance, this design has an important advantage: the qualification output remains essentially interpretable and auditable, which is preferable 
to black-box inference in high-stakes settings \cite{rudin2019stop}.

The objective of this study was to develop and internally validate this governance-aware certification framework for biospecimen qualification. 
The proposed framework separates critical governance admissibility from biological-operational burden, preserves SPREC-coded pre-analytical 
evidence, records the resulting qualification as a local certification layer, and routes specimens into five score-based certification grades: 
Grade A for optimal profiles, Grade B for mild suboptimal profiles, Grade C for moderate suboptimal profiles, Grade D for severe suboptimal 
profiles, and Grade E for critical or governance-failure conditions.

\section{Methods}\label{sec:methods}

\subsection{Model development}\label{subsec:model_development}

This methodological study was designed to develop and internally validate a deterministic score-based certification framework for biospecimen 
quality management in translational biobank settings. The model was developed through a standards-informed knowledge-engineering process rather 
than through a statistical training procedure. SPREC, MIABIS, ISO 20387, and ISBER Best Practices were not treated as direct scoring systems. 
Instead, they were used as complementary evidence sources for structuring a deterministic, auditable, and interpretable classification framework 
\cite{betsou2010,lehmann2012,norlin2012,merinomartinez2016,campbell2018,iso20387}. In this model, SPREC provides the standardized 
pre-analytical evidence axes, while the proposed Ribeiro Classification Score (RCS) converts matrix-specific SPREC-derived axes into 
a weighted penalty burden using heuristic priors informed by literature and standards, and subsequently maps the resulting score 
into five certification grades ranging from A to E.

This separation was necessary because the standards used in the model have different methodological roles and controlled 
terminology. The metadata-governance admissibility layer derived from MIABIS, ISO 20387 and ISBER Best Practices is detailed in Supplementary 
Methods Section~S1.

\subsection{Score-Based Stratification Model}\label{subsec:category_stratification}

The proposed model is a score-based certification model rather than a direct categorical classifier. Its mathematical structure 
separates three components: a non-compensable governance gate, a SPREC-informed weighted penalty structure, and a five-level A--E 
certification function. The classification score generated by the proposed model is hereafter referred to as the RCS.

The SPREC was used to define the matrix-specific pre-analytical axes of the model, but the codes were not treated as universal 
numerical scores. Instead, each derived axis received a maximum weight according to its operational criticality, informed by 
biospecimen science literature indicating that collection, processing, ischemia, fixation, temperature, and storage conditions may 
affect molecular composition, biomarker stability, analytical reproducibility, and downstream interpretability 
\cite{moore2011brisq,hassis2015,carithers2019,bass2014}. Each observed code then received a severity value within its own domain, 
preserving SPREC as a standardized descriptive structure while allowing the RCS to quantify the accumulated operational burden of 
governance-admissible specimens.

The axis-level penalty rationale and matrix-specific maximum weights for fluid and solid biospecimens are provided in Supplementary Methods 
Section~S2. The complete intra-axis severity mapping is provided in Supplementary File~1, Supplementary Table~S3.

\subsection{Computational workflow and governance-aware scoring architecture}

The RCS was designed as a deterministic governance-aware computational framework for biospecimen quality classification, following 
principles of interpretable and rule-based decision systems in biomedical informatics 
\cite{shortliffe2014,russell2021,rudin2019stop}. The workflow separates governance admissibility assessment, matrix-specific penalty 
computation, and final classification into distinct stages.

The process begins with biospecimen metadata and SPREC-derived descriptors, which are evaluated for semantic consistency and 
contextual admissibility. Unresolved, incompatible, invalid, or unsupported descriptors are not interpreted as zero-severity 
conditions and may trigger governance-related restrictions.

After governance assessment, matrix-specific penalty computation is performed using predefined intra-axis severity mappings and 
weighted penalty rules. The resulting biological-operational penalty score ($P_{bio}$) is converted into the Ribeiro Classification 
Score through the deterministic relationship:

\begin{equation}
RCS = 100 - P_{bio}
\label{eq:rcs_definition}
\end{equation}

The final classification is assigned according to predefined threshold intervals corresponding to Grades A--E. Governance failures 
remain non-compensable and override score-based classification.

\subsection{Ribeiro Classification Score (RCS)}\label{subsec:rcs_model}

For a given biospecimen matrix \(k\), where \(k\) denotes either the fluid or solid pathway, the penalty assigned to each SPREC-
derived axis is defined as:

\begin{equation}
p_i^{(k)} = W_i^{(k)} \cdot s_i(x_i),
\label{eq:axis_penalty}
\end{equation}

where \(p_i^{(k)}\) is the penalty assigned to axis \(i\), \(W_i^{(k)}\) is the corresponding maximum axis weight, \(x_i\) is the 
observed SPREC-coded condition, and \(s_i(x_i)\) is the normalized intra-axis severity.

The intra-axis severity function was constrained to the interval:

\begin{equation}
0 \leq s_i(x_i) \leq 1.
\label{eq:severity_interval}
\end{equation}

This interval normalizes intra-axis severity. For governance-admissible specimens, the biological-operational penalty is defined as:

\begin{equation}
P_{bio}^{(k)} = \sum_{i=1}^{n_k} W_i^{(k)} \cdot s_i(x_i),
\label{eq:total_penalty}
\end{equation}

where \(P_{bio}^{(k)}\) is the accumulated operational penalty for matrix \(k\), and \(n_k\) is the number of penalty-generating 
axes evaluated for that matrix.

The Ribeiro Classification Score is then computed as:

\begin{equation}
RCS^{(k)} = 100 - P_{bio}^{(k)}.
\label{eq:rcs}
\end{equation}

In Equation~\ref{eq:rcs}, \(RCS^{(k)}\) denotes the matrix-specific score for fluid or solid biospecimens, whereas \(P_{bio}^{(k)}\) 
represents the accumulated biological-matrix operational penalty. Therefore, the RCS quantifies accumulated pre-analytical deviation 
from a reference condition rather than directly estimating molecular integrity.

For fluid biospecimens, the RCS is computed from the fluid-specific penalty components as:

\begin{equation}
RCS^{(fluid)} = 100 - (P_{pre}+P_{cent1}+P_{cent2}+P_{post}+P_{store}),
\label{eq:rcs_fluid}
\end{equation}

For solid biospecimens, the RCS is computed from the solid-specific penalty components as:

\begin{equation}
RCS^{(solid)} = 100 - (P_{warm}+P_{cold}+P_{fix}+P_{fixTime}+P_{store}),
\label{eq:rcs_solid}
\end{equation}

The metadata-governance layer was modeled as a non-compensable admissibility gate. Let \(g_j\) denote the binary status of 
governance criterion \(j\):

\begin{equation}
g_j =
\begin{cases}
1, & \text{if governance criterion } j \text{ is satisfied} \\
0, & \text{otherwise}.
\end{cases}
\label{eq:governance_criterion}
\end{equation}

The governance gate is then defined as:

\begin{equation}
G_{gov} = \prod_{j=1}^{q} g_j,
\label{eq:governance_product}
\end{equation}

If any critical criterion fails, \(G_{gov}=0\), and the specimen is routed to Grade E regardless of its score. Thus, Grade E may 
result either from critical penalty accumulation (\(RCS<50\)) or governance failure (\(G_{gov}=0\)).

\begin{equation}
Q_{final}^{(k)} =
\begin{cases}
\text{Grade E}, & \text{if } G_{gov}=0 \\
\operatorname{ScoreGrade}(RCS^{(k)}), & \text{if } G_{gov}=1,
\end{cases}
\label{eq:final_grade_rcs}
\end{equation}

where \(Q_{final}^{(k)}\) is the final certification grade, \(G_{gov}\) is the governance-gate status, and 
\(\operatorname{ScoreGrade}(RCS^{(k)})\) denotes the score-based A--E classification.

\begin{table}[!htbp]
\caption{A--E grading structure for RCS certification.}
\label{tab:rcs_grade_scale}
\centering
\footnotesize
\setlength{\tabcolsep}{5pt}
\renewcommand{\arraystretch}{1.14}

\begin{tabular*}{0.86\textwidth}{@{\extracolsep{\fill}}%
>{\centering\arraybackslash}p{1.8cm}
>{\centering\arraybackslash}p{3.8cm}
>{\raggedright\arraybackslash}p{4.6cm}@{}}
\toprule
\textbf{Grade} 
& \textbf{RCS interval} 
& \textbf{Interpretation} \\
\midrule
Grade A & \(RCS \geq 90\) & Optimal profile \\
Grade B & \(80 \leq RCS < 90\) & Mild suboptimal profile \\
Grade C & \(65 \leq RCS < 80\) & Moderate suboptimal profile \\
Grade D & \(50 \leq RCS < 65\) & Severe suboptimal profile \\
Grade E & \(RCS < 50\) & Critical penalty-burden profile \\
\botrule
\end{tabular*}

\vspace{0.35em}
\begin{minipage}{0.86\textwidth}
\footnotesize
\textit{Table note.} Each score interval corresponds directly to a range of accumulated matrix-specific operational penalties.
\end{minipage}
\end{table}

\subsection{Validation Strategy}\label{subsec:validation_strategy}

The validation strategy was designed to evaluate the logical consistency, A--E score stratification, governance sensitivity, 
threshold behavior, internal calibration behavior, sensitivity to heuristic priors, and computational feasibility of the RCS. 
Because this study did not use real patient-derived biobank records, all experiments were conducted using synthetic biospecimen 
profiles. This strategy was selected because the RCS is a deterministic weighted-penalty model rather than a statistical model 
trained from empirical outcomes; therefore, its validation focused on conceptual consistency, operational behavior, and model 
verification rather than empirical parameter learning \cite{sargent2013}.

The validation was structured around the components defined in the previous sections. The operational priority of each SPREC-derived axis was defined by the maximum contribution \(W_i^{(k)}\) described in Supplementary Methods Table~S2. The calculation variables used in the fluid- and solid-specific RCS equations were described in Supplementary Methods Table~S3. Thus, the validation strategy did not redefine the RCS weights or severity rules; it evaluated whether the predefined model generated the expected certification behavior.

Synthetic biospecimen profiles were generated for fluid and solid matrices according to six predefined validation scenarios summarized in Supplementary Methods Table~S4. These scenarios were designed to evaluate score-derived behavior, governance sensitivity, and A--E certification consistency across the predefined validation conditions.

For each biospecimen matrix \(k\), the admissible combination set was defined as:

\begin{equation}
C_k = X_1^{(k)} \times X_2^{(k)} \times \cdots \times X_{n_k}^{(k)},
\label{eq:combination_set}
\end{equation}

where \(C_k\) is the set of admissible SPREC-derived severity combinations for matrix \(k\) to verify whether the deterministic score function populated the A--E certification scale as expected.

Threshold calibration assessment was performed to verify whether the predefined RCS intervals generated the expected A--E certification behavior. The internal calibration index was computed 
as:

\begin{equation}
I_{cal} = \frac{1}{N}\sum_{r=1}^{N} \mathbb{I}(Q_r = Q_r^{expected}),
\label{eq:calibration_index}
\end{equation}

where \(I_{cal}\) is the internal calibration index, \(N\) is the number of synthetic biospecimen profiles, \(Q_r\) is the grade assigned by the model to profile \(r\), \(Q_r^{expected}\) is 
the expected grade according to the synthetic validation scenario, and \(\mathbb{I}(\cdot)\) is the indicator function, equal to 1 when the assigned grade matches the expected grade and 0 

otherwise. This index was interpreted as a measure of internal consistency between predefined validation scenarios and model-derived certification outputs, not as an empirical calibration 
metric.

Threshold-margin analysis was used to quantify how close each synthetic profile was to a grade boundary. Let \(T=\{10,20,35,50\}\) be the set of accumulated-penalty thresholds corresponding 
to the A/B, B/C, C/D, and D/E transitions. The threshold margin for profile \(r\) was defined as:

\begin{equation}
M_r = \min_{\tau \in T} |P_{bio,r}^{(k)}-\tau|,
\label{eq:threshold_margin}
\end{equation}

where \(M_r\) is the minimum distance between the accumulated matrix-specific penalty of profile \(r\) and the nearest certification 
threshold. Smaller values indicate profiles located near a grade boundary and therefore more sensitive to threshold transitions. Additional threshold-transition outputs are provided in Supplementary Results Section~S5.

Governance sensitivity was evaluated using controlled admissibility-failure scenarios affecting metadata completeness, semantic admissibility, traceability, monitoring, and 
documentation. 

\subsection{Runtime Benchmark Analysis}\label{runtime_benchmak}

Runtime benchmarking evaluated computational scalability across cohorts ranging from \(N=10^4\) to \(5\times10^6\) biospecimens. Therefore, 
the local decision logic was expected to remain constant with respect to the total database size, whereas batch processing was expected to 
scale linearly with the number of biospecimens processed.

The implementation was organized as a modular reproducible validation pipeline composed of governance-validation routines, matrix-specific 
scoring functions, synthetic-profile generators, threshold-analysis modules, and figure-generation scripts. The computational workflow was 
designed to support reproducible execution and future integration into biobank information systems and metadata-management environments. 
Additional details regarding the computational environment, benchmark configuration, package versions, and reproducibility parameters are 
provided in Supplementary Methods Table~S5.

\section{Results}\label{sec:results}

The validation results are organized around the conceptual structure of the final RCS model illustrated in Figure~\ref{fig:rcs_workflow}. Admissible biospecimen profiles were evaluated through matrix-specific SPREC-informed penalty components and assigned to threshold-based grades from A to E, while controlled admissibility-failure scenarios were evaluated separately. 
Taken together, the results support the internal consistency, threshold interpretability, robustness, and computational feasibility of the RCS.

\begin{figure}[htbp!]
\justifying
\includegraphics[height=0.65\textheight]{Figure_1_Conceptual_RCS_Workflow.pdf}
\caption{Conceptual computational workflow of the Ribeiro Classification Score (RCS). Biospecimen metadata and SPREC-derived pre-analytical 
descriptors undergo semantic validation and governance admissibility assessment before matrix-specific penalty computation. Admissible records 
are processed through fluid- or solid-specific biological quality models, whereas non-admissible records are assigned Grade E through the 
governance gate. Matrix-specific penalties are aggregated into biological quality scores and combined through deterministic score computation 
to generate the final RCS value and its corresponding certification output (Grades A--E).}
\footnotesize{
\textit{Figure note.}
For fluid biospecimens, penalties include pre-centrifugation delay
($P_{pre}$), primary centrifugation ($P_{cent1}$), secondary centrifugation
($P_{cent2}$), post-centrifugation delay ($P_{post}$), and storage
($P_{store}$). For solid biospecimens, penalties include warm ischemia
($P_{warm}$), cold ischemia ($P_{cold}$), preservation method ($P_{fix}$),
fixation time ($P_{fixTime}$), and storage ($P_{store}$). These components
are aggregated into matrix-specific biological quality scores
($P_{bio}^{(fluid)}$ and $P_{bio}^{(solid)}$), as detailed in Supplementary
File 2, Supplementary Methods Tables S2 and S3.}
\label{fig:rcs_workflow}
\footnotesize
\end{figure}

\subsection{Logical verification and synthetic validation}

The property-based verification layer passed all predefined logical checks summarized in Table~\ref{tab:validation_summary}. These checks 
confirmed that governance failure routes specimens directly to Grade E, that the RCS is computed as \(100-P_{bio}\), that penalties remain 
bounded within the expected interval, that severity behaves monotonically, that identical inputs produce identical outputs, that fluid and 
solid matrices use distinct matrix-specific axes, and that not-scored or unresolved conditions are not interpreted as zero-severity 
conditions. 

Instead, such conditions are routed through governance or contextual admissibility logic. Complete logical-verification outputs and 
supplementary validation details are provided in Supplementary Results Sections~S1 and~S2.

\begin{table}[!htbp]
\caption{Summary of validation evidence for the final RCS model.}
\label{tab:validation_summary}

\footnotesize
\setlength{\tabcolsep}{4.5pt}
\renewcommand{\arraystretch}{1.22}

\begin{tabular}{@{}p{3.4cm}p{3.7cm}p{5.4cm}@{}}

\toprule
\textbf{Validation domain}
& \textbf{Main result}
& \textbf{Interpretation} \\
\midrule

Logical verification
& 7/7 checks passed
& Confirmed deterministic behaviour, including governance override, bounded penalties, monotonic severity response, reproducibility, and 
matrix-specific axis selection. \\

Synthetic validation
& 6000 profiles evaluated
& Tested predefined validation scenarios across fluid and solid biospecimen matrices. \\

Internal calibration
& \(I_{\mathrm{cal}}=0.988833\)
& Showed high agreement between expected scenario grade and assigned RCS grade, with 5933 matches and 67 mismatches. \\

Governance sensitivity
& 1000 governance-failure profiles
& All controlled admissibility-failure profiles were assigned to Grade E. \\

Critical penalty burden
& 1003 critical penalty profiles
& Confirmed that severe accumulated degradation is represented within the score-derived component of the model. \\

\bottomrule

\end{tabular}

\vspace{0.35em}

\noindent\footnotesize{
\textit{Table note.}
Detailed outputs for the corresponding validation domains are provided in Supplementary Results Sections~S1 and~S2.
}
\end{table}

Synthetic validation showed strong agreement between expected scenario labels and assigned RCS grades. The internal calibration index was 0.988833, indicating that the deterministic score 
structure reproduced the intended A--E stratification in nearly all synthetic profiles. The remaining mismatches occurred near predefined score boundaries and were therefore consistent with 
deterministic threshold-margin behavior rather than with random instability. The distribution of expected scenarios against assigned grades is shown in 
Figure~\ref{fig:synthetic_validation_heatmap}.

\begin{figure}[!htbp]
\centering
\includegraphics[width=\textwidth]{Figure_2_Synthetic_Validation_Heatmap.pdf}
\caption{
Synthetic validation of the RCS.
Heatmap showing the proportion of synthetic validation profiles assigned to each final grade across predefined validation scenarios for fluid and solid biospecimens.
}

\label{fig:synthetic_validation_heatmap}
\end{figure}

\subsection{Sensitivity analysis and threshold stability}

Governance sensitivity analysis confirmed that all controlled admissibility-failure scenarios were assigned to Grade E (Supplementary Results Section~S3).

The observed influence pattern followed the predefined axis weights. At 100\% isolated severity, the most influential fluid axis was pre-centrifugation delay, which produced \(P_{bio}=30\),
\(RCS=70\), and Grade C. Fluid storage also reached Grade C at maximal isolated severity, whereas primary centrifugation, post-centrifugation delay, and second centrifugation produced lower 
isolated-grade effects. For solid biospecimens, warm and cold ischemia each produced Grade C at maximal isolated severity, while fixation or stabilization type, fixation time, and storage 
produced Grade B when evaluated in isolation. Complete axis-level influence and global sensitivity outputs are provided in Supplementary Results Section~S4.

Weight-perturbation robustness showed minimal grade instability across 2000 perturbation iterations, with 1000 iterations per matrix. Fluid profiles had a mean changed-rate of 0 and a mean 
absolute RCS shift of 0.046060. Solid profiles had a mean changed-rate of 0.000982 and a mean absolute RCS shift of 0.045161. These values indicate that moderate perturbations of heuristic 
priors changed numerical RCS values slightly, but rarely changed final grade assignment. Detailed perturbation-robustness outputs are provided in Supplementary Results Section~S6.

Threshold-stability analysis showed limited sensitivity to moderate weight perturbations and preserved the expected transition structure across fluid and solid matrices (Supplementary 
Results S5).

\subsection{Model behaviour and computational scalability}

The component-level analysis summarized in Table~\ref{tab:ablation_summary} evaluated whether the final RCS grade distributions were driven primarily by the non-compensable governance gate, 
by the axis weights, or by their combined structure. Component-level analysis showed that removing the governance gate substantially altered grade distributions, whereas replacing the original weights with equal weights produced only minor changes (Table 3).

\begin{table}[!htbp]
\caption{Component-level contribution to final RCS grade distributions.}
\label{tab:ablation_summary}
\centering

\footnotesize
\setlength{\tabcolsep}{4pt}
\renewcommand{\arraystretch}{1.12}

\begin{tabular}{@{}p{4.2cm}ccccc@{}}
\toprule
\textbf{Model configuration}
& \textbf{Grade A}
& \textbf{Grade B}
& \textbf{Grade C}
& \textbf{Grade D}
& \textbf{Grade E} \\
\midrule

Full governance-aware RCS
& 407 & 392 & 404 & 396 & 801 \\

Without governance gate
& 806 & 393 & 404 & 396 & 401 \\

Equal axis weights
& 405 & 393 & 405 & 396 & 801 \\

\bottomrule
\end{tabular}

\vspace{0.35em}

\begin{minipage}{0.92\textwidth}
\footnotesize
\textit{Table note.}
Values represent the number of synthetic profiles assigned to each final grade under each model configuration. The ``without governance gate'' configuration removes the non-compensable governance override, whereas the equal-axis-weight configuration preserves the governance gate but replaces the original axis weights with equal weights.
\end{minipage}
\end{table}

\begin{figure}[!htbp]
\centering
\includegraphics[
    width=\textwidth,
    keepaspectratio
]{Figure_3_Combinatorial_Distribution_Admissible_Grades.pdf}
\caption{Combinatorial distribution of admissible RCS grades. The figure shows the proportion and number of admissible combinatorial profiles assigned to each RCS grade for 
fluid and solid biospecimens. The distribution represents the theoretical admissible combination space and should not be interpreted as expected real-world biobank 
frequencies.}
\label{fig:combinatorial_distribution}
\end{figure}

The distribution of admissible combinatorial profiles across RCS grades is shown in Figure~\ref{fig:combinatorial_distribution}. Runtime benchmarking across \(10^4\) to
\(5\times10^6\) biospecimens demonstrated near-linear scalability, with stable per-profile computational cost and an excellent linear fit (\(R^2=0.998758\); adjusted
\(R^2=0.998721\)). Detailed benchmark outputs are provided in Supplementary Results Section~S7.

Taken together, the validation results support the internal logical consistency of the RCS, the interpretability of threshold-based A--E stratification, robustness to moderate perturbation of heuristic priors, and computational feasibility for large synthetic cohorts.

Deterministic equivalence between sequential and multithreaded execution was confirmed across all benchmark workloads (Supplementary Results Section~S8).

\section{Discussion}\label{sec:discussion}

This study developed and internally validated a governance-aware heuristic penalty model for biospecimen quality management in biobanks. The central 
contribution of the proposed RCS is not to replace SPREC, MIABIS, ISO 20387, or ISBER Best Practices, but to create an explicit and auditable score-
based decision layer over these standards. SPREC provides structured pre-analytical evidence, MIABIS supports interoperable metadata representation, 
and ISO 20387 and ISBER Best Practices support governance, traceability, documentation, monitoring, and quality-management requirements. The RCS uses 
these complementary standards to separate two distinct questions: whether a specimen is admissible for score computation, and how much penalty burden 
is accumulated among admissible specimens. The RCS should therefore be interpreted as an operational qualification framework, not as a direct 
surrogate marker of molecular integrity, clinical validity, or downstream assay performance.

The validation results support this separation. Property-based checks confirmed that the mathematical definition of the RCS was implemented 
deterministically, that penalties remained bounded, and that matrix-specific axes were correctly selected. The high internal calibration index 
indicates that the predefined A--E thresholds behaved as intended across synthetic validation scenarios. Importantly, this calibration index should not be interpreted as external empirical accuracy. It measures internal agreement between predefined validation scenarios and deterministic model outputs, which is appropriate for a rule-based system developed through knowledge engineering rather than statistical training.

The governance-failure analysis is one of the most important findings of the study. In all controlled governance-failure scenarios, the final grade was Grade E even when the classification 
grade would have been Grade A. This behavior is methodologically intentional. A specimen with incomplete profile may not be safely interpreted through an additive penalty score. The model 
therefore prevents a high numerical RCS from compensating for failures in the evidentiary basis required for score computation.

The combinatorial and threshold-margin analyses show that the RCS grade structure is interpretable as a deterministic partition of accumulated 
penalty burden. The theoretical admissible combination space was concentrated in higher-burden categories because the analysis enumerated all 
admissible combinations, including combinations that may be uncommon in operational biobanks. Therefore, the Grade D and Grade E proportions 
in the combinatorial analysis should not be interpreted as expected real-world frequencies. Instead, they demonstrate that the predefined 
penalty structure can populate the full A--E scale and that threshold transitions occur according to the cumulative burden encoded in the 
model.

Sensitivity results further support the internal coherence of the weighting structure. The global sensitivity ranking followed the intended 
priority of the axes. In fluid biospecimens, pre-centrifugation delay and storage were dominant contributors, whereas in solid biospecimens 
warm and cold ischemia were the strongest isolated contributors. These findings are expected from the predefined weights and provide a model-
verification result. The weight-perturbation analysis showed that moderate perturbations produced small numerical RCS shifts and minimal 
changes in final grade assignment. This suggests that the A--E certification structure is not excessively fragile to small changes in 
heuristic priors.

The revised threshold-stability matrix provides a transparent view of how axis-weight perturbations affect grade-transition points. This is 
important because unreached transitions are also informative: they indicate that isolated low- or moderate-weight axes cannot by themselves 
push the 
score across severe boundaries within the 0--100\% severity range.

The analysis confirms that the governance gate is not redundant. This result demonstrates that governance failure is not simply another 
additive penalty; it is a distinct admissibility 
condition. Conversely, the equal-weight ablation produced a distribution close to the full model, indicating that the broad classification 
behavior is stable, although the final weighted 
model remains preferable because its axis priorities preserve the biological-
operational rationale derived from the literature and standards-informed modeling process.

Runtime benchmarking indicates that the RCS is computationally feasible for large-scale biobank software use. The near-linear runtime behavior 
and stable per-profile computational cost reflect the deterministic structure of the algorithm: for each specimen, the model evaluates a fixed 
number of governance and matrix-specific SPREC-derived conditions. Therefore, local decision logic remains constant with respect to total 
database size, while batch processing scales linearly with the number of biospecimens. This is desirable for implementation in biobank 
information systems, where transparent auditability and predictable computational 
cost are both important. 

The supplementary parallel benchmark confirmed deterministic equivalence between sequential and parallel execution and supports future 
distributed implementations. Nevertheless, the stabilization of execution behavior across workloads up to 5 million biospecimen profiles 
supports the suitability of the framework for large-scale batch partitioning and distributed execution strategies. 

Future implementations using compiled or shared-memory backends may further reduce communication overhead and improve practical multithreaded 
throughput.

\subsection*{Limitations and future directions}

This study has limitations. First, the validation was performed using synthetic profiles rather than real biobank records. This was appropriate for internal verification of a deterministic 
model, but it does not establish external predictive validity for molecular outcomes or downstream assay performance. Second, penalty weights and intra-axis severity mappings were defined 
as structured heuristic priors informed by standards and biospecimen science literature, not as empirically fitted coefficients.

Third, the current framework does not yet evaluate how local biobank staff may differ in assigning workflow context or interpreting borderline documentation. Such human and institutional 
variability may introduce bias in practical implementation. Fourth, although the model supports technical governance and decision auditability, prospective evaluation in real biobank 
workflows will be necessary to assess usability, calibration against operational decisions, and potential impact on specimen selection.

Despite these limitations, the proposed RCS offers a transparent bridge between standards-based biospecimen documentation and operational 
decision support. Its value lies in making explicit 
the rules by which standards-derived evidence is transformed into a local certification status. This approach is particularly relevant for 
translational research environments where 
reproducibility, traceability, and interpretability must be preserved alongside 
computational efficiency.

The deterministic and modular architecture of the RCS may facilitate future integration into laboratory information systems, biobank 
management platforms, and automated governance-monitoring workflows. Also, future studies should evaluate the behavior of the RCS in real-
world biobank infrastructures and investigate potential associations between RCS-derived grades and downstream molecular or analytical 
outcomes.

\section{Conclusions}\label{sec:conclusions}

The RCS provides a deterministic, governance-aware, and standards-informed framework for biospecimen quality certification in biobank 
environments. By combining a non-compensable governance gate with SPREC-derived weighted penalties, the model separates admissibility failures 
from gradual pre-analytical burden and maps the resulting evidence into an interpretable A--E certification scale.

Internal validation demonstrated logical correctness, high agreement between predefined synthetic scenarios and assigned grades, explicit 
governance-failure override, coherent threshold behavior, robustness to moderate weight perturbation, and near-linear runtime scalability. The 
results support the RCS as an auditable decision-support layer that complements, rather than replaces, SPREC, MIABIS, ISO 20387, and ISBER 
Best Practices.

Future work should evaluate the model in real biobank implementations, compare RCS outputs with expert review and molecular quality 
indicators, and assess how human documentation practices influence governance-gate decisions. The current study establishes the technical and 
methodological foundation for such external validation.

\section*{List of abbreviations}

RCS: Ribeiro Classification Score; SPREC: Standard PREanalytical Code; MIABIS: Minimum Information About BIobank data Sharing; ISO: 
International Organization for Standardization; ISBER: International Society for Biological and Environmental Repositories; SOP: Standard 
Operating Procedure; RT: room temperature.

\section*{Declarations}

\section*{Ethics approval and consent to participate}

This study was approved by the Research Ethics Committee of São Francisco University (Universidade São Francisco), 
Brazil, under CAAE 97291926.5.0000.5514 and approval number 8.380.323 granted on April 26, 2026. The Research Ethics 
Committee waived the requirement for written informed consent because the study used exclusively synthetic and simulated 
data and did not involve the collection, access, or analysis of identifiable human participant data. All procedures were 
conducted in accordance with applicable institutional guidelines and ethical regulations.


\section*{Consent for publication}

Not applicable.

\section*{Declaration of generative AI and AI-assisted technologies in the writing process}

During the preparation of this work, the authors used generative AI tools exclusively for language refinement and editorial assistance. The 
authors reviewed and edited the content as needed and take full responsibility for the final content of the publication.

\section*{Data availability}

The synthetic validation outputs, figure-source datasets, supplementary machine-readable tables, runtime benchmark logs, computational-
environment records, and reproducibility scripts are available in the project GitHub repository: \url{https://github.com/deniseribeiro-beep/rcs-validation}. The repository contains only synthetic validation data and does not include real patient-derived biobank records.

\section*{Code availability}

The \textsf{R} scripts used for synthetic profile generation, RCS computation, validation checks, sensitivity analysis, ablation analysis, runtime benchmarking, and figure generation are 
available in the project GitHub repository: \url{https://github.com/deniseribeiro-beep/rcs-validation}.

\section*{Funding}

This study was financed in part by the Coordenação de Aperfeiçoamento de Pessoal de Nível Superior – Brasil (CAPES) – Finance Code 001.

\section*{CRediT authorship contribution statement}

DR conceived the study, organized the methodological framework, implemented the validation strategy, and drafted the manuscript. FA contributed to methodological refinement and computational design. LMC contributed to mathematical design and review algorithmic interpretation. VIBS supervised the study and contributed to conceptual and editorial review. All authors 
read and approved the final manuscript.

\section*{Acknowledgements}

Not applicable.

\section*{Supplementary material}

\noindent\textbf{Supplementary File 1. SPREC reference and severity tables.}
This supplementary document contains the SPREC 2.0 reference descriptors for fluid and solid biospecimens and the intra-axis severity mapping used as the standards-based 
foundation for the RCS framework.

\medskip

\noindent\textbf{Supplementary File 2. Supplementary Methods.}
This supplementary document contains the detailed methodological framework, including metadata-governance admissibility definitions, matrix-specific penalty structures, 
RCS calculation variables, certification algorithms, synthetic validation scenarios, and computational-environment specifications.

\medskip

\noindent\textbf{Supplementary File 3. Supplementary Results.}
This supplementary document contains the complete supplementary validation outputs, including logical verification analyses, post-run validation checks, synthetic 
validation and internal calibration summaries, governance-sensitivity analyses, threshold-margin analyses, axis-influence and global-sensitivity outputs, threshold-
transition stability analyses, weight-perturbation robustness results, runtime scalability benchmarks, persistent-cluster parallel benchmark outputs, and the source 
tables supporting the supplementary results.

\bibliographystyle{elsarticle-num}
\bibliography{references}

\end{document}
